% Mise en page des PDF : LaTeX classique (Latin Modern), A4, compacte,
% entièrement lisible en noir et blanc. Chargé après le préambule de Quarto.
\usepackage{xcolor}
\definecolor{vertprofond}{HTML}{1C322D}
\definecolor{grisfilet}{HTML}{A1A1AA}

% Paragraphes sans alinéa, séparés par un petit espace.
\setlength{\parindent}{0pt}
\setlength{\parskip}{0.5ex plus 0.2ex minus 0.1ex}

% Titres resserrés, en vert profond (assez sombre pour l'impression N&B).
\setkomafont{disposition}{\normalfont\bfseries\color{vertprofond}}
\setkomafont{section}{\large}
\setkomafont{subsection}{\normalsize}
\setkomafont{subsubsection}{\normalsize\itshape}
\RedeclareSectionCommand[beforeskip=-1.6ex plus -0.4ex minus -0.2ex,
  afterskip=0.6ex plus 0.1ex]{section}
\RedeclareSectionCommand[beforeskip=-1.2ex plus -0.3ex minus -0.2ex,
  afterskip=0.4ex plus 0.1ex]{subsection}
\RedeclareSectionCommand[beforeskip=-1ex plus -0.2ex,
  afterskip=0.3ex]{subsubsection}

% Listes serrées.
\usepackage{enumitem}
\setlist{topsep=0.3ex, itemsep=0.15ex, parsep=0pt, partopsep=0pt}

% Figures et tableaux : peu d'espace autour, images plafonnées en hauteur
% pour qu'un schéma n'occupe jamais une demi-page.
\setlength{\textfloatsep}{1.2ex plus 0.4ex minus 0.4ex}
\setlength{\floatsep}{1ex plus 0.4ex minus 0.4ex}
\setlength{\intextsep}{1ex plus 0.4ex minus 0.4ex}
\setlength{\abovecaptionskip}{0.4ex}
\setlength{\belowcaptionskip}{0pt}
\setkomafont{caption}{\small\itshape}
\setkomafont{captionlabel}{\small\itshape}
\setkeys{Gin}{width=\linewidth, height=0.28\textheight, keepaspectratio}
\renewcommand{\arraystretch}{1.15}

% En-tête et pied de page : formation et chapitre en haut, titre et
% numéro de page en bas. Les valeurs viennent du partial title.tex.
\usepackage{lastpage}
\usepackage[headsepline=0.4pt, footsepline=0.4pt]{scrlayer-scrpage}
\clearpairofpagestyles
\setkomafont{pageheadfoot}{\footnotesize\color{vertprofond}}
\setkomafont{pagenumber}{}
\ihead{\vintageformation}
\ohead{\vintagechapitre}
\ifoot*{\vintagetitre}
\ofoot*{\thepage\,/\,\pageref*{LastPage}}
\pagestyle{scrheadings}

% Schémas TikZ : les styles communs suivent dans tikz-vintage.tex.
\usepackage{tikz}
